\documentclass[border=10pt]{standalone}
\usepackage[no-math]{fontspec}
\usepackage[brazil]{babel}
\usepackage{amsmath,amssymb}
\IfFileExists{C:/Windows/Fonts/Lato_Regular.ttf}{%
  \setmainfont{Lato}[
    Path = C:/Windows/Fonts/,
    UprightFont = Lato_Regular.ttf,
    BoldFont = Lato_Bold.ttf,
    ItalicFont = Lato_Italic.ttf,
    BoldItalicFont = Lato_BoldItalic.ttf
  ]
  \setsansfont{Lato}[
    Path = C:/Windows/Fonts/,
    UprightFont = Lato_Regular.ttf,
    BoldFont = Lato_Bold.ttf,
    ItalicFont = Lato_Italic.ttf,
    BoldItalicFont = Lato_BoldItalic.ttf
  ]
}{%
  \RequirePackage[default]{lato}
}
\renewcommand{\familydefault}{\sfdefault}

\usepackage{xcolor}
\usepackage{tikz}
\usepackage{pgfplots}
\pgfplotsset{compat=1.18}
\usetikzlibrary{arrows.meta,positioning,shapes.geometric,fit,backgrounds,calc,decorations.pathreplacing,patterns}

% Paleta institucional UEFS / Bioengenharia de Solos
\definecolor{bsNavy}{HTML}{1F4268}
\definecolor{bsNavyDark}{HTML}{152E4A}
\definecolor{bsGold}{HTML}{B38A3D}
\definecolor{bsGoldDark}{HTML}{8C6927}
\definecolor{bsTeal}{HTML}{0E7C7B}
\definecolor{bsForest}{HTML}{2D6A4F}
\definecolor{bsAlert}{HTML}{9E2A2B}
\definecolor{bsEarth}{HTML}{8C5A3C}
\definecolor{bsWater}{HTML}{2E86AB}
\definecolor{bsInk}{HTML}{26313D}
\definecolor{bsMuted}{HTML}{5B6573}

\definecolor{bgNavy}{HTML}{EDF2F7}
\definecolor{bgGold}{HTML}{F8F3E9}
\definecolor{bgTeal}{HTML}{EBF6F5}
\definecolor{bgForest}{HTML}{EEF5F0}
\definecolor{bgAlert}{HTML}{FDF2F2}

\tikzset{
  >=Stealth,
  bsBox/.style={
    draw=#1, fill=white, rounded corners=5pt, line width=1.15pt,
    inner sep=7pt, align=center, font=\footnotesize\sffamily, text=bsInk
  },
  bsTint/.style={
    draw=#1, fill=#1!8!white, rounded corners=5pt, line width=1.15pt,
    inner sep=7pt, align=center, font=\footnotesize\sffamily, text=bsInk
  },
  bsBadge/.style={
    fill=#1, text=white, rounded corners=9pt, inner xsep=8pt, inner ysep=3pt,
    font=\scriptsize\bfseries\sffamily
  },
  bsTitle/.style={font=\normalsize\bfseries\sffamily, text=bsNavy, align=center},
  bsSub/.style={font=\footnotesize\sffamily, text=bsMuted, align=center},
  bsArrow/.style={->, >=Stealth, line width=1.35pt},
  bsDArrow/.style={<->, >=Stealth, line width=1.35pt}
}


\begin{document}
\begin{tikzpicture}[node distance=0.6cm]

  % Título e subtítulo
  \node[bsTitle] (title) at (0, 3.85)
    {Arquitetura da Bioengenharia de Solos Tropicais e Conservação de Águas};
  \node[bsSub] (sub) at (0, 3.35)
    {Integração entre pedologia, hidrologia, geotecnia e soluções baseadas na natureza};

  % Pilar 1: Solo Tropical & Pedologia
  \node[bsTint=bsEarth, text width=4.6cm, minimum height=3.15cm, anchor=north] (p1) at (-5.3, 2.55) {
    \vspace{0.35cm}
    \textbf{\color{bsEarth}Solo Tropical \& Pedologia}\\[5pt]
    \scriptsize
    \textbullet\ Latossolos, Argissolos e Cambissolos\\[2pt]
    \textbullet\ Intemperismo químico intenso e\\
    \hspace{0.25cm}agregação por óxidos (Fe/Al)\\[2pt]
    \textbullet\ Física do solo: textura, porosidade,\\
    \hspace{0.25cm}infiltração e erodibilidade ($K$)
  };
  \node[bsBadge=bsEarth] at (p1.north) {1. BASE PEDOLÓGICA};

  % Pilar 2: Hidrologia & Mecânica
  \node[bsTint=bsWater, text width=4.6cm, minimum height=3.15cm, anchor=north] (p2) at (0, 2.55) {
    \vspace{0.35cm}
    \textbf{\color{bsWater}Hidrologia \& Mecânica}\\[5pt]
    \scriptsize
    \textbullet\ Balanço hídrico e escoamento\\
    \hspace{0.25cm}superficial (Manning, $n$ e $V$)\\[2pt]
    \textbullet\ Critério de Mohr-Coulomb e\\
    \hspace{0.25cm}tensão cisalhante ($\tau_0$)\\[2pt]
    \textbullet\ Reforço pela coesão radicular ($\Delta c_r$)
  };
  \node[bsBadge=bsWater] at (p2.north) {2. PROCESSOS FÍSICOS};

  % Pilar 3: Intervenções Integradas
  \node[bsTint=bsForest, text width=4.6cm, minimum height=3.15cm, anchor=north] (p3) at (5.3, 2.55) {
    \vspace{0.35cm}
    \textbf{\color{bsForest}Intervenções Integradas}\\[5pt]
    \scriptsize
    \textbullet\ Controle de escoamento: terraços\\
    \hspace{0.25cm}e canaletas vegetadas (bioswales)\\[2pt]
    \textbullet\ Estabilização de encostas e ravinas:\\
    \hspace{0.25cm}paliçadas, fascinas e biomantas\\[2pt]
    \textbullet\ Contenção viva: gabião vivo e Krainer
  };
  \node[bsBadge=bsForest] at (p3.north) {3. ENGENHARIA ECOLÓGICA};

  % Convergência
  \coordinate (c1) at (-5.3, -1.15);
  \coordinate (c2) at (0, -1.15);
  \coordinate (c3) at (5.3, -1.15);
  \coordinate (hub) at (0, -1.55);

  \draw[bsArrow, bsEarth] (p1.south) -- (c1) -| ($(hub) + (-2.0, 0)$);
  \draw[bsArrow, bsWater] (p2.south) -- (hub);
  \draw[bsArrow, bsForest] (p3.south) -- (c3) -| ($(hub) + (2.0, 0)$);

  % Síntese unificada
  \node[bsTint=bsNavy, text width=15.0cm, minimum height=1.35cm, anchor=north, fill=bgNavy] (synth) at (0, -1.6) {
    \vspace{0.18cm}
    \textbf{\color{bsNavy}\small Síntese Aplicada: Projeto Integrador de Controle Erosivo e Conservação Hídrica}\\[4pt]
    \scriptsize
    Diagnóstico multicritério do sítio \quad$\cdot$\quad Dimensionamento hidráulico-geotécnico \quad$\cdot$\quad Restauração funcional e infraestrutura verde (SbN)
  };
  \node[bsBadge=bsNavy] at (synth.north) {CONVERGÊNCIA DA DISCIPLINA};

\end{tikzpicture}
\end{document}
